% Part: incompleteness
% Chapter: representability-in-q
% Section: representable-comp

\documentclass[../../../include/open-logic-section]{subfiles}

\begin{document}

\olfileid{inc}{req}{rpc}
\olsection{الدوال القابلة للتمثيل في $\Th{Q}$ قابلة للحساب}

نبدأ بإثبات أن قابلية الدالة للتمثيل في~$\Th{Q}$ تستلزم قابليتها للحساب.
ونقدّم لذلك لمة في الدوال القابلة للتمثيل في~$\Th{Q}$.

\begin{lem}\ollabel{lem:rep-q}
  لتكن $f(x_0,\dots,x_k)$ قابلة للتمثيل في~$\Th{Q}$. حينئذ توجد
  !!a{formula}~$!A_f(x_0,\dots,x_k,y)$ تحقق:
  \[
  \Th{Q} \Proves !A_f(\num{n_0}, \dots, \num{n_k}, \num{m})
  \quad\text{إذا وفقط إذا}\quad m = f(n_0, \dots, n_k).
  \]
\end{lem}

\begin{proof}
أما اتجاه «إذا» فمقتضى
\olref[int]{defn:representable-fn}\olref[int]{defn:rep:a}.
وأما اتجاه «فقط إذا»، فنبرهنه بالخلف. لنفرض
$\Th{Q}\Proves !A_f(\num{n_0},\dots,\num{n_k},\num{m})$ مع
$m\neq f(n_0,\dots,n_k)$، ولنضع $l=f(n_0,\dots,n_k)$.
فمن \olref[int]{defn:representable-fn}\olref[int]{defn:rep:a} نحصل على
$\Th{Q}\Proves !A_f(\num{n_0},\dots,\num{n_k},\num{l})$.
ومن \olref[int]{defn:representable-fn}\olref[int]{defn:rep:b} نحصل على
$\Th{Q}\Proves
\lforall[y][(!A_f(\num{n_0},\dots,\num{n_k},y)\lif\num{l}=y)]$.
فإذا استعملنا قواعد المنطق وفرضنا
$\Th{Q}\Proves !A_f(\num{n_0},\dots,\num{n_k},\num{m})$، لزم
$\Th{Q}\Proves\eq[\num{l}][\num{m}]$.
لكن \olref[bre]{lem:q-proves-neq} يعطينا أيضًا
$\Th{Q}\Proves\eq/[\num{l}][\num{m}]$؛ فتكون $\Th{Q}$ غير متسقة.
وهذا ممتنع: فإن النموذج القياسي يحقق $\Th{Q}$
(انظر \olref[int][def]{def:standard-model})، أي إن $\Sat{N}{\Th{Q}}$،
وكل نظرية قابلة للتحقيق متسقة بمبرهنة السلامة
(\tagrefs{prfAX/{fol:axd:sou:cor:consistency-soundness},
prfSC/{fol:seq:sou:cor:consistency-soundness},
prfND/{fol:ntd:sou:cor:consistency-soundness},
prfTab/{fol:tab:sou:cor:consistency-soundness}}).
\end{proof}

\begin{lem}
إذا كانت الدالة قابلة للتمثيل في $\Th{Q}$ فهي قابلة للحساب.
\end{lem}

\begin{proof}
نبيّن أولًا وجه القول على سبيل التقريب. فحساب~$f$ يجري بالبحث الآتي:
نسرد كل ما يمكن أن يكون من !!{derivation}s~$\delta$ في لغة الحساب،
وسردها ممكن آليًا. ثم ننظر في كل واحدة: أهي !!a{derivation}
لـ!!a{formula} على الصورة
$!A_f(\num{n_0},\dots,\num{n_k},\num{m})$، أعني الـ!!{formula} الممثلة
لـ$f$ في~$\Th{Q}$ بحسب \olref{lem:rep-q}؟ فإن كانت كذلك ثبت
$m=f(n_0,\dots,n_k)$ باللمة \olref{lem:rep-q}، فحصلت قيمة~$f$.
ولا يستمر البحث بلا نهاية؛ إذ
$\Th{Q}\Proves !A_f(\num{n_0},\dots,\num{n_k},
\num{f(n_0,\dots,n_k)})$، فلا بد أن يظهر في السرد $\delta$ المطلوبة.

ولم تبلغ هذه الحجة بعدُ غاية الدقة؛ فالإجراء الموصوف يتناول
!!{derivation}s و!!{formula}s، لا الأعداد وحدها، ولم نفصّل وجه إمكان
«سرد جميع الـ!!{derivation}s الممكنة» آليًا. غير أن الطريق إلى ضبطها
قد تقدم: فالحدود و!!{formula}s و!!{derivation}s ترمّز بأعداد غودل،
والدوال العودية العامة تعطينا نموذجًا مضبوطًا للحساب. وقد ثبت أن العلاقة
$\Prf[\Th{Q}](d,y)$ عودية، بل عودية بدائية؛ ومعناها أن $d$ عدد غودل
لـ!!a{derivation} من بديهيات~$\Th{Q}$ للـ!!{formula} التي عدد غودل لها~$y$.
ونحتاج معها إلى الدالتين العوديتين البدائيتين
$\fn{num}$ (\olref[art][trm]{prop:num-primrec}) و$\fn{Subst}$
(\olref[art][sub]{prop:subst-primrec}).
ومن هذه المعطيات نعرّف~$f$ بالتصغير، فتثبت عودية $f$.

نضع أولًا
\begin{multline*}
  A(n_0, \dots, n_k, m) = \\
  \fn{Subst}(\fn{Subst}(\dots\fn{Subst}(\Gn{!A_f}, \fn{num}(n_0), \Gn{x_0}),\\ \dots),
  \fn{num}(n_k),  \Gn{x_k}), \fn{num}(m), \Gn{y})
\end{multline*}
وليس وراء تعقيد هذه الكتابة إلا الدالة
$A(n_0,\dots,n_k,m)=
\Gn{!A_f(\num{n_0},\dots,\num{n_k},\num{m})}$.

ثم نأخذ العلاقة~$R(n_0,\dots,n_k,s)$، ومعناها أن $(s)_0$
عدد غودل لـ!!a{derivation} من~$\Th{Q}$ للصيغة
$!A_f(\num{n_0},\dots,\num{n_k},\num{(s)_1})$:
\[
R(n_0, \dots, n_k, s) \quad\text{إذا وفقط إذا}\quad \Prf[\Th{Q}]((s)_0, A(n_0,
\dots, n_k, (s)_1))
\]
فمتى وجدنا~$s$ تحقق $R(n_0,\dots,n_k,s)$، وجدنا زوج العددين
$(s)_0$ و$(s)_1$؛ أولهما، $(s)_0$، عدد غودل
  لـ!!a{derivation} للصيغة المعبّر عنها اختصارًا بـ$!A_f(\num{n_0},\dots,\num{n_k},\num{(s)_1})$.
فالبحث عن~$s$ هو نظير البحث عن $d$ و$m$ معًا في الحجة غير الصورية.
وهذا البحث عن $s$ تحدده دالة قابلة للحساب بالتّصغير المنتظم.
وانتظام $R$ بيّن: لكل $n_0$، \dots، $n_k$، توجد
!!a{derivation}~$\delta$ تثبت ما تقوله
$\Th{Q}\Proves !A_f(\num{n_0},\dots,\num{n_k},
\num{f(n_0,\dots,n_k)})$.
فتصدق $R(n_0,\dots,n_k,s)$ عند
$s=\tuple{\Gn{\delta},f(n_0,\dots,n_k)}$، ونحصل على $f$ بالتعريف:
\[
  f(n_0,\dots,n_{k}) = (\umin{s}{R(n_0, \dots, n_k, s)})_1.
  \]
  % تصحيح مصدر (C291-01): وُحّد اسم الصيغة الممثلة وأُثبت واسم الصيغة وغلاف العدد في المثال المختصر.
\end{proof}

\end{document}
